\begin{papercard}{citron2021}
\cf{Stages} \stagetag{H} (typology of privacy harms; severity and aggregation), \stagetag{S} (definition of harm versus risk versus problem), \stagetag{R} (how to articulate harms). A 70-page law-review article (Boston University Law Review 102:793, 2022; SSRN 2021) by Danielle Keats Citron and Daniel J. Solove. Legal in character and US-centred, it is the most developed \emph{typology} of privacy harms available in the corpus and is the intended backbone for the thesis' harm tree. Its legal sections (standing, causes of action) are context; the typology (Part V) and the three structural challenges (Part II) are what the thesis can reuse.

\cf{Problem and thesis} US privacy enforcement is impeded by a harm requirement: ``the requirement of harm has significantly impeded the enforcement of privacy law'' and courts add a harm element even where statutes do not require one (e.g., standing after \emph{Spokeo} and \emph{TransUnion}) \pc{citron2021}{1}{The requirement of harm has significantly impeded the enforcement of privacy law}. Courts struggle because privacy harms often involve future uses of data, are typically small (frustration, anxiety, inconvenience) and dispersed, and ``when these minor harms are suffered at a vast scale, they produce significant harm to individuals, groups, and society'' \pc{citron2021}{1}{When these minor harms are suffered at a vast scale, they produce significant harm to individuals, groups, and society}. Two contributions: a typology of harms, and an approach to \emph{when} harm should be required at all. The intro states that courts ``refuse to recognize privacy harms that do not involve tangible financial or physical injury'' \pc{citron2021}{4}{Courts, however, refuse to recognize privacy harms that do not involve tangible financial or physical injury}.

\cf{Definition} Harms ``involve injuries, setbacks, losses, or impairments to well-being'' and leave people or society worse off than before; a \emph{cognizable} harm is one the law accepts as suitable for intervention \pc{citron2021}{7}{Harms involve injuries, setbacks, losses, or impairments to well-being}. Harms are a subset of Solove's broader ``privacy problems'' (undesirable states of affairs); the harm typology differs from his 2008 taxonomy of problems (footnotes 28, 203).

\cf{Part I: where law demands harm} Standing (Article III injury in fact: concrete, particularized, actual or imminent: \emph{Clapper} 2013, \emph{Spokeo} 2016, \emph{TransUnion} 2021), causes of action (contract, tort, statutes with or without statutory damages) and regulatory enforcement (FTC unfairness: a small harm to many may be substantial). Illustrative rulings are summarised for context (Doe v.\ Chao, FAA v.\ Cooper, Senne). Not needed in detail for the thesis.

\cf{Part II: three structural difficulties} (1) \textbf{Aggregation of small harms}: ``privacy harms often involve the aggregation of many small harms to each individual, which is compounded by the aggregation of all these harms to many individuals'' (the single unwanted email versus hundreds; tiny harm times millions) \pc{citron2021}{24}{Privacy harms often involve the aggregation of many small harms to each individual, which is compounded by the aggregation of all these harms to many individuals}; class actions are an imperfect remedy. (2) \textbf{Risk: unknowable and future harms}: harms depend on future uses by third parties; ``risk is involved with many different types of privacy harm'' \pc{citron2021}{25}{Risk is involved with many different types of privacy harm}; harms are ``highly contextual'', depending on the data, its use and combination with other data \pc{citron2021}{26}{Privacy harms are highly contextual, with the harm depending upon how the data is used, what data is involved, and how the data might be combined with other data}. (3) \textbf{Individual vs.\ societal harms}: privacy is constitutive of society; individualistic doctrine omits societal impact \pc{citron2021}{26}{Privacy harms often involve injury not just to individuals but to society}.

\cf{Part III: align goals and remedies} Enforcement has three goals: compensation (damages for harm), deterrence (prevent future violations), equity (right wrongs by other means). Using a single remedy (damages) for all goals is the root problem: ``a wrench ... to hammer in a nail'' \pc{citron2021}{27}{Privacy law enforcement has three predominant goals}. When deterrence or equity is the goal, harm should not be required. Private rights of action act as ``private attorneys general''. Part IV argues that proper identification of the interest at stake matters because courts accept ill-fitting \emph{tangible} harms (battery drain, lost minutes) as stand-ins: ``the law perversely redresses trivial setbacks while ignoring major problems'' \pc{citron2021}{36}{The law perversely redresses trivial setbacks while ignoring major problems}, and recognising harm has expressive value because it tells people that ``their suffering is real'' \pc{citron2021}{36}{Law’s recognition of privacy harms tells individuals that their suffering is real}.

\cf{Part V: the typology} Seven types, with subtypes for psychological and autonomy harms; ``privacy harms not only differ in type but also in their severity'' \pc{citron2021}{38}{Privacy harms not only differ in type but also in their severity}. Figure 1 groups them \pc{citron2021}{39}{Our typology groups privacy harms into seven basic types}:
\begin{enumerate}[nosep,leftmargin=*]
\item \textbf{Physical harms}: bodily injury or death; well recognised. Data disclosure enables stalking and violence (home address from motor-vehicle records; \emph{Remsburg v.\ Docusearch}). Doxing is both a risk of physical harm and psychological harm (fear) \pc{citron2021}{42}{Doxing—the disclosure of personal data to facilitate people being located, contacted, and harassed—creates a serious threat of physical harm}.
\item \textbf{Economic harms}: monetary loss or loss of value (identity theft, fraud); courts accept direct losses, but not time and anxiety spent mitigating, nor loss of opportunities and productivity \pc{citron2021}{42}{Economic harms involve monetary losses or a loss in the value of something}, \pc{citron2021}{43}{Many privacy violations involve the loss of important opportunities rather than direct financial injuries}.
\item \textbf{Reputational harms}: injuries to reputation and standing; lost business or employment or social rejection \pc{citron2021}{45}{Reputational harms involve injuries to an individual’s reputation and standing in the community}.
\item \textbf{Psychological harms} \pc{citron2021}{49}{Psychological harms involve a range of negative mental responses, such as anxiety, anguish, concern, irritation, disruption, or aggravation}, in two subtypes: \emph{emotional distress} (painful feelings: annoyance, frustration, fear, embarrassment, anger, anxiety) \pc{citron2021}{49}{Emotional distress encompasses a wide range of emotions, including annoyance, frustration, fear, embarrassment, anger, and various degrees of anxiety} and \emph{disturbance} (unwanted intrusions that disturb tranquility, interrupt activities and sap time, such as unsolicited calls, texts and spam) \pc{citron2021}{52}{Disturbance involves unwanted intrusions that disturb tranquility, interrupt activities, sap time, and otherwise serve as a nuisance}. Courts accept emotional distress in the privacy torts but are sceptical elsewhere.
\item \textbf{Autonomy harms}: ``restricting, undermining, inhibiting, or unduly influencing people's choices''; people are denied the freedom to decide or ``tricked into thinking that they are freely making choices'' \pc{citron2021}{53}{People are either directly denied the freedom to decide or are tricked into thinking that they are freely making choices when they are not}. Six subtypes: (a) \emph{coercion}, constraint or undue pressure on freedom to act or choose (e.g., conditioning treatment on data uses) \pc{citron2021}{54}{Coercion involves a constraint or undue pressure on one's freedom to act or choose}; (b) \emph{manipulation}, undue influence over behaviour or decisions, ``one of the most prevalent forms of autonomy harm in the consumer privacy context''; distinguishing persuasion and nudging from manipulation is ``challenging and contestable'' \pc{citron2021}{54}{Distinguishing between acceptable influencing (persuasion and nudging) and unacceptable influencing (manipulation) is challenging and contestable}; Spencer: ``an intentional attempt to influence a subject's behavior by exploiting a bias or vulnerability'' \pc{citron2021}{55}{an intentional attempt to influence a subject’s behavior by exploiting a bias or vulnerability}; coercion is blunt and noticed, manipulation ``subtle and sneaky'' \pc{citron2021}{55}{Coercion is blunt and forthright: one almost always knows one is being coerced. Manipulation is subtle and sneaky}; (c) \emph{failure to inform}, not giving individuals information to make informed choices about data or rights \pc{citron2021}{56}{Failure to inform involves failing to provide individuals with information to assist them in making informed choices about their personal data or exercise of their privacy rights}; (d) \emph{thwarted expectations}, breaking promises about collection, use and disclosure \pc{citron2021}{57}{The harm caused by thwarted expectations involves the undermining of people’s choices, such as breaking promises made about the collection, use, and disclosure of personal data}, noting that privacy policy plays a small role in forming expectations whereas settings, advertising and design elements shape them \pc{citron2021}{59}{Website or browser privacy settings, company advertising, statements, and other design elements have an influence on people’s expectations}; (e) \emph{lack of control}, inability to make choices about one's data or to curtail uses \pc{citron2021}{61}{Lack of control involves the inability to make certain choices about one’s personal data or to be able to curtail certain uses of the data}; (f) \emph{chilling effects}, inhibition of lawful activities such as speech, association or exploring ideas, illustrated by Grindr users who would stop sharing HIV status after it was shared with analytics firms \pc{citron2021}{63}{the gay dating app Grindr had shared subscribers’ HIV status with analytics firms}.
\item \textbf{Discrimination harms}: entrenching inequality and disadvantaging people by gender, race, sexual orientation, age or group membership; differs from autonomy harms by shame, stigma and societal entrenchment; e.g., employers and insurers accessing period-tracking data \pc{citron2021}{63}{Discrimination harms involve entrenching inequality and disadvantaging people based on gender, race, national origin, sexual orientation, age, group membership}.
\item \textbf{Relationship harms}: damage to relationships important to health and functioning, personal, professional and with organisations; ``twofold: the loss of confidentiality and in the longer term, damage to the trust'' (example: fertility-tracking app communities that rely on confidentiality) \pc{citron2021}{67}{Relationship harms involve the damage to relationships that are important for one’s health, well-being, life activities, and functioning in society}.
\end{enumerate}

\cf{Conclusion} Focus enforcement on deterrence and widespread compliance, with compensation for significant harm \pc{citron2021}{69}{Its primary focus should be on the deterrence of violations with the goal of encouraging widespread compliance}; the common law, strict liability, restitution and fiduciary duties offer paths.

\cf{Limits (card)} A normative legal argument, not empirical; US doctrine (Article III, Spokeo) dominates and does not carry over to EU law (GDPR Art.\ 82 allows non-material damage); the typology is a classification for courts, not a validated measurement instrument, and the authors do not test reliability or show that the seven types are mutually exclusive (they note overlaps, e.g.\ doxing = physical + psychological; discrimination vs.\ autonomy); severity is mentioned but not scaled; no probability model for risk.

\cf{Use for this thesis} (1) Take the seven top-level types and the six autonomy subtypes as a \emph{candidate} skeleton for the stage-H harm tree, attaching each harm node to a concrete app feature, data flow or screen finding (e.g., health data shared to analytics $\to$ chilling effect, discrimination; consent screen with default-on toggles $\to$ manipulation, lack of control; broken in-app promise $\to$ thwarted expectations; vague disclosure $\to$ failure to inform). (2) Use the three structural challenges to argue why per-person severity scores are inappropriate: aggregation, risk and contextuality justify a qualitative, evidence-linked tree rather than a probability. (3) Record in the Decision Log that harm categories are adapted from a legal typology to a health-app setting and are not claims of legal cognizability; the typology's overlaps should be handled by allowing multi-labels. (4) Pair with Solove's taxonomy, Martin's harm work and Nissenbaum's contextual integrity for the \emph{why} of each node, and with the app-review/walkthrough evidence for the \emph{where}.
\end{papercard}
